% ID: FIS-FLU-VAS-001
% Titolo: Livello del liquido nei vasi comunicanti
% Disciplina: Fisica
% Area: Fluidostatica
% Argomento: Vasi comunicanti
% Sottoargomento: Equilibrio idrostatico
% Classe: Terza liceo artistico
% Difficolta: 2
% Tipo: quesito_teorico
% Risultato: soluzione da aggiungere
% Tempo_stimato: 6 min
% Competenze: interpretazione fisica, argomentazione, equilibrio idrostatico
% Prerequisiti: pressione idrostatica, equilibrio, liquidi
% Tag: fisica, fluidostatica, vasi_comunicanti, pressione_idrostatica, equilibrio
% Fonte: verifica_fisica_fluidostatica_anonimizzata
% Autore: Riccardo Carli
% Licenza: CC-BY-SA-4.0

\begin{esercizio}
Quattro recipienti di forme diverse sono collegati tra loro alla base e
contengono lo stesso liquido. Dopo un tempo sufficiente, il liquido risulta alla
stessa altezza in tutti i recipienti.

Spiega perche il livello finale del liquido e identico nei quattro contenitori,
anche se le loro forme sono diverse.
\end{esercizio}

\begin{soluzione}
% Soluzione da aggiungere.
\end{soluzione}

